\chapter{Running Scenarios}
łabel{ch:running−scenarios}
∈dex{GM guidance!scenarios}∈dex{scenarios!running}

A scenario is a self−contained unit of play with a clear stake, meaningful choices, and visible pressure. Your goal as GM is not to control the outcome, but to \textbf{shape the pressure} so that player decisions create memorable consequences. This chapter distills Fate's Edge into actionable GM techniques for facilitating scenarios—from preparation to execution—using the system's core tools: Story Beats, timers, and player agency.

\section{GM Mindset: Co−Creating Drama}
Before rolling dice, internalize these principles. They transform rule enforcement into collaborative storytelling.

\begin{tcolorbox}[colback=purple!5!white, colframe=purple!75!black, title=GM Core Principles]
\begin{itemize}
  ℹtem \textbf{Narrative primacy:} Mechanics serve the story. When in doubt, choose what creates the most dramatic consequence or interesting choice.
  ℹtem \textbf{Pressure over punishment:} Use Story Beats to escalate tension (e.g., “The ritual advances”) not to penalize players (“You lose 2 HP”).
  ℹtem \textbf{Fail forward:} Every roll changes the situation. A Miss isn’t “nothing happens”—it’s “the alarm sounds and guards converge.”
  ℹtem \textbf{Player agency first:} Offer costs, timers, or complications instead of saying “no.” (“You can try, but it will take time—tick the Ritual Timer +1.”)
  ℹtem \textbf{Your preparation is invisible magic:} Between−session work turns random encounters into meaningful episodes. Players never see the strings—only the music.
\end{itemize}
\end{tcolorbox}

\section{Preparing a Scenario: The 15−Minute GM Sprint}
Effective scenario prep focuses on \textbf{pressure points}, not plots. Do this before each session:

\begin{enumerate}
  ℹtem \textbf{Identify the stake:} What is the players' \textit{immediate} goal? (e.g., “Retrieve the Sunstone before the ritual completes.”) Write it as a question: “Can they get the Sunstone in time?”
  ℹtem \textbf{Set 1−−3 visible timers:} Use the Three−Timer Rule (see \ref{sec:three−timer}) to track progress/opposition. Label them clearly: \emph{Ritual Completion [6]}, \emph{Guard Patrols [4]}.
  ℹtem \textbf{Prepare 2−−3 flexible situations:} Not scripts, but pressure cookers with clear stakes and timers. Example:
    \begin{itemize}
      ℹtem \textit{Situation:} “The Sunstone is on a pedestal in the collapsing chamber. The ceiling cracks ominously.”
      ℹtem \textit{Timer:} \emph{Chamber Collapse [8]} (ticks on Partial/Miss inside chamber)
      ℹtem \textit{Stakes:} Succeed = get Sunstone; Fail = take Harm 1 + timer advances.
    \end{itemize}
  ℹtem \textbf{Seed player hooks:} Review Bonds/Complications. Ask: “How does this situation threaten what they care about?” (e.g., A player bonded to the Scholar Mira? The rival Captain Voss taunts them about her fate.)
  ℹtem \textbf{Check complication debt:} Note unresolved Complications from last session—they add +1 SB to the first scene.
\end{enumerate}

\begin{tcolorbox}[colback=blue!5!white, colframe=blue!75!black, title=Scenario Prep Checklist]
  ∈dex{preparation!scenario}
  \begin{itemize}
    ℹtem [ ] Stake phrased as a yes/no question
    ℹtem [ ] 1−−3 visible timers named and sized (4/6/8/10)
    ℹtem [ ] 2−−3 flexible situations (stake + timer)
    ℹtem [ ] At least one player hook per PC
    ℹtem [ ] Complication debt noted (adds SB to Scene 1)
  \end{itemize}
\end{tcolorbox}

\section{Running the Scene: The 90−Second Loop}
Every meaningful player action follows this GM procedure. \textit{Internalize this loop—it’s your engine.}

\begin{enumerate}
  ℹtem \textbf{Declare intent and approach:} “What do you want to do, and how?” (Player states Attribute+Skill)
  ℹtem \textbf{Set Position and DV:} 
    \begin{itemize}
      ℹtem \textbf{Position:} Dominant (advantage), Controlled (normal), Desperate (high risk). Default: Controlled.
      ℹtem \textbf{DV:} 2 (routine) to 5+ (extreme). Default: 3 (pressured). \textit{Never inflate DV to “challenge”—use Position instead.}
    \end{itemize}
    Announce: “Controlled, DV 3. On a Miss, the guard hears you.”
  ℹtem \textbf{Roll and count:} Pool = Attribute+Skill d10s. Successes: 6–9 = 1, 10 = 2. Count 1s for Story Beats (SB).
  ℹtem \textbf{Apply Outcome Matrix:} 
    \begin{itemize}
      ℹtem Successes $≥$ DV, SB=0: Clean Success—intent achieved, no complication.
      ℹtem Successes $≥$ DV, SB>$>$0: Success with SB—intent achieved; GM spends SB.
      ℹtem 0 $<$ Successes $<$ DV: Partial—progress; player gains 1 Boon.
      ℹtem Successes = 0: Miss—no progress; player gains 2 Boons; GM spends SB.
    \end{itemize}
    \textit{Remember: A Partial moves the story forward. A Miss always escalates.}
  ℹtem \textbf{Spend Story Beats:} Use the SB Spend Menu (see \ref{sec:sb−menu}) to introduce complications that \textit{change the situation}. Tie spends to visible outcomes: “I spend 2 SB—the alarm blares and two guards rush in (tick Guard Patrols +2).”
  ℹtem \textbf{Describe the new situation:} End with one sentence stating the new fiction: “The chamber shudders as dust rains from the ceiling—Ritual Completion is now at 5 of 6.”
\end{enumerate}

\begin{tcolorbox}[colback=green!5!white, colframe=green!75!black, title=Between−Session Checklist]
  ∈dex{preparation!between−sessions}
  \begin{enumerate}
    ℹtem \textbf{Update timers:} Tick faction, travel, or campaign timers based on last session’s outcomes.
    ℹtem \textbf{Check complication debt:} Note unresolved player Complications; they add +1 SB to the next session’s first scene.
    ℹtem \textbf{Prepare 2−−3 flexible situations:} Not scripts, but situations with a clear stake and a timer.
    ℹtem \textbf{Review player goals:} Look at Bonds, Patron obligations, and personal arcs. Seed at least one hook per player.
    ℹtem \textbf{Refresh your memory:} Re−read the session summary from last time (ask a player to recap if you didn’t write one).
  \end{enumerate}
\end{tcolorbox}

\section{Key GM Techniques}
These are the levers you pull to shape pressure. Practice them until they’re instinctive.

\subsection{The Three−Timer Rule}
łabel{sec:three−timer}
At any given time, maintain \textbf{at most three active timers} in a scene. More creates cognitive overload and dilutes tension. Name them and tick them \textit{aloud} where players can see.

\begin{enumerate}
  ℹtem \textbf{Situation Timer:} Tracks progress toward the scene’s immediate goal (e.g., \emph{Chamber Collapse [8]}, \emph{Persuade the Duke [6]}).
  ℹtem \textbf{Counter−Clock:} Tracks opposition pressure (e.g., \emph{Guard Patrols [4]}, \emph{Ritual Completion [6]}).
  ℹtem \textbf{Campaign Timer:} Tracks long−term front pressure (e.g., \emph{Cult Influence [6]}, \emph{Sun’s Fading [10]}).
\end{enumerate}
Merge redundant timers (e.g., two “guard arrival” timers become one). Retire timers that resolve naturally. \textit{Example:} In a negotiation scene:
\begin{itemize}
  ℹtem Situation Timer: \emph{Reach Agreement [6]} (ticks on Partial in Persuade rolls)
  ℹtem Counter−Clock: \emph{Duke’s Patience [4]} (ticks when players ignore social cues)
  ℹtem Campaign Timer: \emph{Noble Feud Spreads [8]} (ticks between sessions based on faction actions)
\end{itemize}

\subsection{Story Beats Spend Menu}
łabel{sec:sb−menu}
Spend SB to introduce complications that \textit{escalate pressure or reveal truth}. Match SB cost to narrative impact—never spend more than needed.

\begin{center}
\begin{tabular}{p{1.5cm}p{10cm}}
⊤rule
\textbf{SB Cost} & \textbf{Effect} \\
∣rule
1 & Minor pressure: Noise, distraction, tick a 4−−6 segment timer +1, worsen Position by one step. \\
& \textit{Example:} “Your foot scrapes a loose stone—the guard looks up (tick Guard Alert +1).” \\
2 & Moderate setback: Alarm raised, lose cover/advantage, lesser foe appears, split the party briefly. \\
& \textit{Example:} “The guard calls for backup—two more arrive (new timer: Reinforcements [6]).” \\
3+ & Serious trouble+: Reinforcements arrive, key gear breaks, major twist shifts the scene, authority arrives, environment changes dramatically. \\
& \textit{Example:} “The chamber ceiling groans—a 10−ton block begins to fall (Chamber Collapse now at 7 of 8).” \\
⊥tomrule
\end{tabular}
\end{center}
\textit{Pro tip:} One strong beat (e.g., 3 SB for reinforcements) creates more drama than three weak beats (1+1+1 SB for noise). Be generous—unused SB is wasted momentum.

\subsection{Handling Player Creativity}
When players propose unexpected solutions:
\begin{itemize}
  ℹtem \textbf{Say “Yes, but” not “no”:} Offer a cost, timer, or complication. (“You can try to bribe the guard—but it will take time. Make a Persuade roll, and on a Miss, the guard alerts his captain.”)
  ℹtem \textbf{Use Setup/Assist actions:} Let non−specialists contribute meaningfully. (“You can’t pick the lock, but you can hold the torch steady—grant +1 die to the rogue’s roll.”)
  ℹtem \textbf{Embrace the unexpected:} If they bypass your planned encounter, let it work—but escalate pressure elsewhere. (“You trick the guards into leaving their post… but the ritual chamber’s seals begin to weaken. Tick Ritual Completion +1.”)
  ℹtem \textbf{Fail forward aggressively:} A Miss on their creative solution shouldn’t stop the story—it should change direction. (“You fail to convince the guard… but in your panic, you notice a loose stone in the wall—a possible escape route.”)
\end{itemize}

\subsection{Visible Pressure: Making Timers Matter}
Timers only create tension if players \textit{see and understand} them.
\begin{itemize}
  ℹtem \textbf{Name timers descriptively:} \emph{“Ritual Completion”} not \emph{“Timer 3”}. Players should know what advancing it means.
  ℹtem \textbf{Tick them aloud:} “The cult’s chant grows louder—Ritual Completion advances to 4 of 6.”
  ℹtem \textbf{Link timer advances to fiction:} Never just say “tick timer.” Describe what changes: “As the timer hits 5, the air smells of ozone and sulfur—the demon’s presence warps reality nearby.”
  ℹtem \textbf{Place trackers centrally:} Use physical tokens (see \ref{ch:campaigns−timers} for tracker ideas) where all players see them. Update them \textit{during} play, not just between scenes.
\end{itemize}

\section{Common Pitfalls and Fixes}
Even experienced GMs fall into these traps. Recognize and correct them instantly.

\begin{tcolorbox}[colback=red!5!white, colframe=red!75!black, title=Pitfall Fixes]
\begin{itemize}
  ℹtem \textbf{Pitfall:} “I’m hoarding SB for a big climax.” \\
  \textbf{Fix:} Spend SB early and often. Pressure creates drama—saving it makes scenes flat. Unspent SB beyond your limit (Tier +1) carries over, but hoarding starves the table of tension.
  ℹtem \textbf{Pitfall:} “The players are stuck—they keep failing the same roll.” \\
  \textbf{Fix:} Ask: “What is your \textit{goal}?” Then name the simplest obstacle. If they want to sneak past guards but keep failing Stealth, offer: “You can’t sneak—but you could create a distraction (Make a Deceive roll).” 
  ℹtem \textbf{Pitfall:} “I don’t know the rule for X.” \\
  \textbf{Fix:} Make a ruling that keeps the story moving, then check later. Default: DV 3, Position Controlled. Note it for downtime resolution—never halt play to look up rules.
  ℹtem \textbf{Pitfall:} “My prep went unused—the players went off−track.” \\
  \textbf{Fix:} Your prep was never about specific scenes—it was about \textit{pressure}. If they ignored the cultist ambush, the cult’s timer still advances. Reveal consequences: “You avoided the ambush… but the ritual completes without interference. Tick Cult Influence +2.”
  ℹtem \textbf{Pitfall:} “I feel like I’m arguing with players about rules.” \\
  \textbf{Fix:} Remember: You are co−creators, not adversaries. When rules conflict with fiction, \textit{always} side with fiction. Make a ruling that serves the story, then reconcile with rules text afterward.
\end{itemize}
\end{tcolorbox}

\section{Rookie GM Comfort Dials}
If you’re new to Fate’s Edge, start here to build confidence.

\begin{tcolorbox}[colback=orange!5!white, colframe=orange!75!black, title=Rookie GM Settings]
\begin{itemize}
  ℹtem \textbf{Soft SB:} For first 2 sessions, cap SB spend per roll at 1–2 (unless it’s a set−piece moment like a timer filling). Prevents overwhelming new players.
  ℹtem \textbf{Visible Timers:} Put \textit{all} active timers on the table. Name them aloud. Tick them in ink so players see the progression.
  ℹtem \textbf{Tag Cards:} Print one−liners for frequent tags (e.g., \texttt{[WARD]: Blocks outsiders. DV = target’s Cap.}) and hand them out when relevant.
  ℹtem \textbf{One Move, One Sentence:} After every ruling, end with one sentence stating the new situation: “The gate holds—but the timer for Cult Influence advances to 3 of 6 as the chant peaks.”
\end{itemize}
\end{tcolorbox}

\section{Scenario Types: Adapting the Loop}
Apply the 90−Second Loop to these common scenario frames with these tweaks:

\subsection{Social Encounters}
Use Persuasion Timers (see SRD 17.3.3):
\begin{itemize}
  ℹtem Clean Success: +2 segments to Persuasion Timer
  ℹtem Partial: +1 segment; GM may spend SB for complication (“They agree—but you’ll owe a favor”)
  ℹtem Miss: No progress; GM ticks Opposition Timer (when full, target hardens Position)
  ℹtem \textit{GM Tip:} Spend SB to reveal hidden motives (“The duke’s smile doesn’t reach his eyes—he’s stalling for time”) or escalate stakes (“His hand twitches toward his sword—you see the guard shift behind you”).
\end{itemize}

\subsection{Travel Legs}
Frame each leg as a scenario with roles (see \ref{sec:gm−travel}):
\begin{itemize}
  ℹtem \textbf{Guide:} Wits+Survival to tick Travel Timer +1 (success) or backwards (failure)
  ℹtem \textbf{Scout:} Wits+Stealth to avoid surprise on first encounter
  ℹtem \textbf{Quartermaster:} Wits+Craft to manage Supply
  ℹtem \textbf{Watch:} Presence+Notice for +1 Position on first ambush defense roll
  ℹtem \textit{GM Tip:} On failed role, spend SB to introduce immediate encounter (“As you argue the route, a wolf pack emerges from the fog”).
\end{itemize}

\subsection{Magical Experiments}
Treat spellcasting as pressure (see SRD 17.5.1):
\begin{itemize}
  ℹtem Free casters: Position worsens by 1 when engaged unless spell has [TOUCH] or [WARD]
  ℹtem Rite casters: Mark Obligation when invoking; Push It for +1 Obligation for increased effect
  ℹtem \textit{GM Tip:} On Partial/Miss, spend SB for thematic backlash (“The fireball singes your eyebrows—next action is Desperate”) or environmental shift (“The spell’s resonance cracks a stalactite—it begins to fall”).
\end{itemize}

\section{Ending the Scenario: Leave Hooks, Not Loose Ends}
Close scenarios to create momentum for the next:
\begin{enumerate}
  ℹtem \textbf{Resolve the immediate stake:} State clearly whether they achieved their goal (yes/no/partial).
  ℹtem \textbf{Advance relevant timers:} Especially Campaign and Faction timers based on outcomes.
  ℹtem \textbf{Seed 1−−2 hooks for next session:} 
    \begin{itemize}
      ℹtem “The rival captain escaped—but she dropped this bloodstained sigil.”
      ℹtem “The ritual failed… but the cult’s leader whispers your true name as she flees.”
      ℹtem “You saved the village—but the Sunstone’s light now attracts \textit{something} from the deep dark.”
    \end{itemize}
  ℹtem \textbf{Explicitly note complication debt:} “Your Complication ‘Sealed Hold Refuses You’ remains unresolved—it will add +1 SB to your first action next session.”
  ℹtem \textbf{Ask one reflective question:} “What did you learn about your character tonight?” (Use answers to seed arcs/bonds).
\end{enumerate}

\section{The GM’s Sacred Trust}
Remember: Your off−table work is where the magic happens. Between sessions:
\begin{itemize}
  ℹtem Update timers based on scene outcomes (factions advance, pressures spread)
  ℹtem Check complication debt (unresolved Complications = +1 SB next scene)
  ℹtem Prepare flexible situations (stake + timer, not scripts)
  ℹtem Review player sheets for hooks (Bonds, Patron obligations, personal arcs)
  ℹtem Refresh your memory on last session’s events
\end{itemize}
This invisible preparation transforms dice rolls into consequence−rich stories. When players say, “That session felt \textit{alive},” it’s because you honored this trust—not because of any rule you enforced at the table.

Let the pressure shape the story. Let the dice sing. And never forget: you are playing the game too.